% Appendix C.2: CSD Implementation Details

\subsection{CSD Implementation Details}
\label{app:subsec:csd_implementation}

% 详细描述 Context-Stripped Decoding 的实现：baseline subtraction 操作流程、target prompt 构造方式、
% 以及面向不熟悉 Patchscopes 的读者的背景补充。
% 预期效果：读者无需阅读 Patchscopes 原文即可理解 CSD 的完整操作。

This section provides a self-contained description of the Context-Stripped Decoding (CSD) procedure, so that readers unfamiliar with Patchscopes \citep{ghandeharioun2024patchscopes} can reproduce the method without consulting the original paper.
The formal definition is given in \cref{eq:csd}; here we unpack each step.

\paragraph{Step 1: Target prompt construction.}
Each code-understanding sample consists of a code context $C$ (the program) and a question suffix $Q$ (e.g., \texttt{\# The value of x is "}).
The \emph{source prompt} is the full concatenation $S = [C;\, Q]$.
The \emph{target prompt} retains only the question suffix: $T = Q$, thereby stripping away all code context.
Because $T$ contains no information about the program, any correct answer decoded from $T$ must originate from the patched hidden state rather than from the surrounding context.

\paragraph{Step 2: Patching procedure.}
For a given layer $\ell \in \{0, \dots, L{-}1\}$, CSD proceeds as follows:
\begin{enumerate}
    \item Run the source prompt $S = [C;\, Q]$ through the full model. Extract the hidden state $\mathbf{h}^\ell_S$ at the last token position at layer $\ell$.
    \item Construct a separate forward pass with the target prompt $T = Q$. At layer $\ell$, replace the last-token hidden state in the target run with $\mathbf{h}^\ell_S$ from the source run:
    \begin{equation}
    \label{eq:csd_patch}
    \tilde{\mathbf{h}}^\ell_T \;\leftarrow\; \mathbf{h}^\ell_S.
    \end{equation}
    \item Continue the target forward pass from layer $\ell{+}1$ through layer $L{-}1$, applying the remaining transformer blocks $F_{\ell+1}, \dots, F_{L-1}$, followed by LayerNorm and the unembedding matrix $W_u$. Crucially, the attention context for all subsequent layers is restricted to $T$, not the original $S$.
\end{enumerate}
This yields patched logits $\mathbf{z}_{\mathrm{patch}}^\ell = W_u \circ \mathrm{LN} \circ F_{L-1} \circ \cdots \circ F_{\ell+1}(\tilde{\mathbf{h}}^\ell_T)$.

\paragraph{Step 3: Baseline subtraction.}
Running the target prompt $T$ alone (without any patching) produces a baseline logit vector $\mathbf{z}_b$, which captures the language prior induced by $T$.
For instance, the question suffix \texttt{\# The value of x is "} may assign non-trivial probability to certain digits even without any code context.
To isolate the contribution of the patched hidden state, we subtract this baseline:
\begin{equation}
\label{eq:csd_subtract}
\mathbf{z}_{\mathrm{CSD}}^\ell = \mathbf{z}_{\mathrm{patch}}^\ell - \mathbf{z}_b.
\end{equation}
Finally, we apply softmax and restrict the resulting distribution to the target token set $\mathcal{T} = \mathcal{D} \cup \{\bar{d}\}$, where $\mathcal{D} = \{0, 1, \dots, 9\}$ is the set of digit tokens and $\bar{d}$ is a residual class:
\begin{equation}
\label{eq:csd_distribution}
\Phi_{\mathrm{C}}^\ell(\mathbf{h}^\ell)[t] = \frac{\exp(\mathbf{z}_{\mathrm{CSD}}^\ell[t])}{\sum_{t' \in \mathcal{T}} \exp(\mathbf{z}_{\mathrm{CSD}}^\ell[t'])}, \quad t \in \mathcal{T}.
\end{equation}
This recovers the CSD distribution defined in \cref{eq:csd}.

\paragraph{Step 4: Practical considerations.}
Several implementation details are worth noting:

\begin{itemize}
    \item \textbf{Layer-wise independence.} Patching is performed at each layer $\ell \in \{0, \dots, L{-}1\}$ independently.
    Each layer produces its own CSD distribution $\Phi_{\mathrm{C}}^\ell$; there is no cross-layer interaction between patching runs.

    \item \textbf{Baseline reuse.} The baseline logit $\mathbf{z}_b$ depends only on the target prompt $T$ and is therefore identical across all layers for a given sample.
    We compute $\mathbf{z}_b$ once and reuse it for all $L$ patching runs, reducing the total number of forward passes from $2L$ to $L{+}1$ per sample.

    \item \textbf{Hook-based implementation.} In practice, we implement the hidden-state replacement using PyTorch forward hooks.
    A hook registered at layer $\ell$ intercepts the target run's hidden state and overwrites the last-token position with $\mathbf{h}^\ell_S$.
    This avoids modifying the model architecture and allows patching at arbitrary layers without rerunning earlier layers of the target pass.

    \item \textbf{Positional encoding.} Because the source prompt $S$ and target prompt $T$ have different lengths, the absolute position index of the last token differs between the two runs.
    Following \citet{ghandeharioun2024patchscopes}, we do not adjust positional encodings when patching; the injected hidden state retains the positional information from the source run.
    In the Qwen2.5-Coder architecture, which uses Rotary Position Embeddings (RoPE), this means the patched state carries the source-side rotary phase.
    We verified empirically that this does not degrade the diagnostic's discriminative power for our single-token prediction task.
\end{itemize}
